\documentclass{article}

\usepackage{amsmath,amssymb}
\usepackage{xurl}
\usepackage[hidelinks]{hyperref}
\setlength{\emergencystretch}{2em}

% Shared commands for notes in docs/paper-gaps/.

\DeclareUnicodeCharacter{2080}{\ifmmode{}_{0}\else\textsubscript{0}\fi}
\DeclareUnicodeCharacter{2081}{\ifmmode{}_{1}\else\textsubscript{1}\fi}
\DeclareUnicodeCharacter{2082}{\ifmmode{}_{2}\else\textsubscript{2}\fi}
\DeclareUnicodeCharacter{2099}{\ifmmode{}_{n}\else\textsubscript{n}\fi}

\newcommand{\Lean}{\textsc{Lean}}
\ExplSyntaxOn
\NewDocumentCommand{\leanid}{m}{
  \ifmmode
    \textsf{\detokenize{#1}}
  \else
    \group_begin:
    \urlstyle{sf}
    \tl_set:Nn \l_tmpa_tl {#1}
    \tl_replace_all:Nnn \l_tmpa_tl {\_} {_}
    \exp_args:NV \path \l_tmpa_tl
    \group_end:
  \fi
}
\ExplSyntaxOff
\providecommand{\tr}{\operatorname{tr}}
\newcommand{\tnleanIssueBase}{https://github.com/LionSR/TNLean/issues/}
\newcommand{\tnleanPrBase}{https://github.com/LionSR/TNLean/pull/}
\newcommand{\tnleanDocBase}{https://sirui-lu.com/TNLean/docs/}
\newcommand{\ghissue}[1]{\href{\tnleanIssueBase#1}{GitHub issue \##1}}
\newcommand{\ghpr}[1]{\href{\tnleanPrBase#1}{PR \##1}}
\newcommand{\doclink}[1]{\href{\tnleanDocBase#1.html}{\path{#1}}}

% Dirac notation and the identity operator, as in blueprint/src/macros/common.tex.
% \providecommand keeps note-local definitions authoritative.
\usepackage{dsfont}
\providecommand{\ket}[1]{|#1\rangle}
\providecommand{\bra}[1]{\langle #1|}
\providecommand{\braket}[2]{\langle #1 | #2 \rangle}
\providecommand{\ketbra}[2]{|#1\rangle\!\langle #2|}
\providecommand{\Id}{\mathds{1}}
\providecommand{\one}{\mathds{1}}
\providecommand{\C}{\mathbb{C}}
\providecommand{\MN}[1]{M_{#1}(\C)}
\providecommand{\MD}{M_D(\C)}

% External referencing for paper sources.  The build helper writes sanitized
% auxiliary files under build/paper-aux/ and excludes duplicated source labels.
\usepackage{xr}

% Import a generated paper auxiliary file with a caller-supplied label prefix.
% The candidate paths support builds from docs/paper-gaps/, the repository root,
% and build/paper-aux/ itself.
\newcommand{\importpaperaux}[2]{%
  \IfFileExists{../../build/paper-aux/#2.aux}{%
    \externaldocument[#1]{../../build/paper-aux/#2}%
  }{%
    \IfFileExists{build/paper-aux/#2.aux}{%
      \externaldocument[#1]{build/paper-aux/#2}%
    }{%
      \IfFileExists{#2.aux}{%
        \externaldocument[#1]{#2}%
      }{}%
    }%
  }%
}

% General external reference macros
\newcommand{\paperref}[2]{\ref{#1-#2}}
\newcommand{\papereqref}[2]{(\ref{#1-#2})}

% CPSV16 source (1606.00608 / MPDO-22-12-17-2.tex) external document and shortcuts
\importpaperaux{cpsv16-}{1606.00608/MPDO-22-12-17-2-tnlean}
\newcommand{\cpsvref}[1]{\paperref{cpsv16}{#1}}
\newcommand{\cpsveqref}[1]{\papereqref{cpsv16}{#1}}

% Verdict marker: \gapnote{<kind>}{<status>}. Kind is the policy
% classification (clarification, local-correction, scope-restriction,
% unfaithful, false-source, open-gap); status is open, wip (elimination
% underway), resolved, or historical. Machine-read by the site index;
% prints a verdict line.
\newcommand{\gapnote}[2]{\par\noindent\textbf{Verdict:} #1 (#2).\par}


\title{CPSV16 Vertical-Sector Invertibility and Sector Relabelling}
\author{}
\date{2026-07-19}

\begin{document}
\maketitle
\gapnote{scope-restriction}{resolved}

\section{The source argument}

The proof of the general mixed-state renormalization fixed-point theorem in
CPSV16 uses the finite-dimensional $C^*$-algebras
\begin{align}
    \mathcal A_1
        &= \bigoplus_\alpha M_{d^1_\alpha}(\mathbb C),
    \label{eq:cpsv16_vertical_sector_invertibility_first_sector_algebra}\\
    \mathcal A_2
        &= \bigoplus_\beta M_{d^2_\beta}(\mathbb C).
    \label{eq:cpsv16_vertical_sector_invertibility_second_sector_algebra}
\end{align}
Let $T$ and $S$ denote the physical maps, and let $R_1,R_2$ and
$\widetilde R_1,\widetilde R_2$ denote the sector realization and compression
maps used in the source. The transported maps are
\begin{align}
    \widetilde T
        &= \widetilde R_2 T R_1,
    \label{eq:cpsv16_vertical_sector_invertibility_transported_T}\\
    \widetilde S
        &= \widetilde R_1 S R_2.
    \label{eq:cpsv16_vertical_sector_invertibility_transported_S}
\end{align}
For each horizontal bond matrix $X$, CPSV16 proves
\begin{align}
    \widetilde T\!\left(\bigoplus_\alpha m_\alpha M_\alpha(X)\right)
        &= \bigoplus_\beta n_\beta A_\beta(X),
    \label{eq:cpsv16_vertical_sector_invertibility_T_contractions}\\
    \widetilde S\!\left(\bigoplus_\beta n_\beta A_\beta(X)\right)
        &= \bigoplus_\alpha m_\alpha M_\alpha(X).
    \label{eq:cpsv16_vertical_sector_invertibility_S_contractions}
\end{align}
Equations~\ref{eq:cpsv16_vertical_sector_invertibility_T_contractions}
and~\ref{eq:cpsv16_vertical_sector_invertibility_S_contractions} imply that
the one-site contraction family lies in
$\operatorname{Fix}(\widetilde S\widetilde T)$ and the two-site contraction
family lies in $\operatorname{Fix}(\widetilde T\widetilde S)$
\cite[Appendix~C.4, lines~1972--1980]{Cirac2016MPDO_arXiv}.

The next step does not use the linear span of the contraction family. CPSV16
invokes the density-weighted fixed-point form of Wolf's
Theorem~6.14~\cite{Wolf2012Quantum} and writes the fixed-point space as
\begin{align}
    \mathcal F
        &= U\left(0_r\oplus\bigoplus_k
            \rho_k\otimes M_{h_k}(\mathbb C)\right)U^*.
    \label{eq:cpsv16_vertical_sector_invertibility_fixed_point_form}
\end{align}
For a contraction family $V$, let $C_L(V)$ denote the span of all length-$L$
products from $V$. The basis-of-normal-tensors hypothesis gives some $L>0$
such that
\begin{align}
    C_L(V)
        &= \mathcal A_1.
    \label{eq:cpsv16_vertical_sector_invertibility_contraction_product_span}
\end{align}
Equation~\ref{eq:cpsv16_vertical_sector_invertibility_fixed_point_form} gives
\begin{align}
    C_L(\mathcal F)
        &= U\left(0_r\oplus\bigoplus_k
            \rho_k^L\otimes M_{h_k}(\mathbb C)\right)U^*.
    \label{eq:cpsv16_vertical_sector_invertibility_fixed_product_span}
\end{align}
Each $\rho_k$ is positive definite, so
\begin{align}
    \dim C_L(\mathcal F)
        &= \dim\mathcal F.
    \label{eq:cpsv16_vertical_sector_invertibility_fixed_product_dimension}
\end{align}
The inclusions
\begin{align}
    \mathcal A_1
        = C_L(V)
        \subseteq C_L(\mathcal F)
        \subseteq \mathcal A_1
    \label{eq:cpsv16_vertical_sector_invertibility_product_span_inclusions}
\end{align}
therefore imply $\mathcal F=\mathcal A_1$ and hence
\begin{align}
    \widetilde S\widetilde T
        &= \operatorname{id}_{\mathcal A_1}.
    \label{eq:cpsv16_vertical_sector_invertibility_source_first_identity}
\end{align}
The symmetric argument proves the other composition identity
\cite[Appendix~C.4, lines~1980--1995]{Cirac2016MPDO_arXiv}.

\section{The formal boundary}

Three source-specific ingredients preceding the fixed-point classification
have been proved:
\begin{enumerate}
    \item the two transported composites fix their respective weighted bond
          contraction families;
    \item for a basis of normal tensors with nonzero weights, products of some
          positive length $L$ from these contractions span the full sector
          algebra;
    \item on the contraction families generated by the source, the final
          compression preserves the active image and therefore preserves total
          trace.
\end{enumerate}
The third result does not by itself prove trace preservation on the whole
sector algebra. The contraction family need not span the algebra linearly;
only products of a fixed positive length are known to span. A linear map that
fixes each $V(X)$ need not fix
$V(X_1)\cdots V(X_L)$. Consequently, neither trace preservation nor the
identity conclusion follows from the generators by linearity. The independent
maximal-support argument below proves whole-sector trace preservation, while
the identity conclusion still uses the fixed-point classification.

Two finite-dimensional final implications are formalized. First, an
endomorphism is the identity if it fixes every length-$L$ contraction product
and those products span the ambient space.\footnote{Formal declaration:
\leanid{MPOTensor.eq_id_of_weightedVerticalBondContractionProducts_fixed}
in \doclink{TNLean/MPS/MPDO/VerticalSectorGeneration}.}
CPSV16 does not assume that the products themselves are fixed.

For the second implication, let $\mathcal F$ be the fixed-point space and let
$C_L(\mathcal F)$ be the span of its length-$L$ products. If the distinguished
contractions are fixed, their length-$L$ products span the ambient sector
algebra, and
\begin{align}
    \dim C_L(\mathcal F)
        &\leq \dim\mathcal F,
    \label{eq:cpsv16_vertical_sector_invertibility_abstract_dimension_bound}
\end{align}
then the endomorphism is the identity.\footnote{Formal declarations:
\leanid{MPOTensor.fixedPointProductSpan} and
\leanid{MPOTensor.eq_id_of_weightedVerticalBondContractions_fixed_of_productSpan_finrank_le}
in \doclink{TNLean/MPS/MPDO/VerticalSectorGeneration}.}
This implication does not require products of the distinguished contractions
to be fixed. It isolates the dimension argument used in the source.

The bound in
Eq.~\ref{eq:cpsv16_vertical_sector_invertibility_abstract_dimension_bound} is
formalized for a positive trace-preserving map on a finite direct sum when its
trace adjoint satisfies the Schwarz inequality and it has a positive-definite
fixed family.\footnote{Formal declaration:
\leanid{MPOTensor.fixedPointProductSpan_finrank_le_of_densityBlocks}
in \doclink{TNLean/MPS/MPDO/VerticalSectorDensityBlocks}.}
This result assumes no tensor hypothesis, multiplicativity, or fixedness of
products. It is nevertheless a full-support corollary: unlike Wolf's
unrestricted Theorem~6.14~\cite{Wolf2012Quantum}, it assumes a
positive-definite fixed family. The tensor argument must obtain this family
from the fixed contractions and their positive-length product span.

The positive trace-preserving mean-ergodic retraction on a full matrix algebra
and its positive unital trace-adjoint retraction onto the adjoint fixed-point
space are also formalized. Their construction does not require a bounded-orbit
hypothesis for the adjoint map. On full support, the adjoint retraction agrees
with the weighted partial-trace block form of Wolf's
Proposition~1.5~\cite{Wolf2012Quantum}. The trace-adjoint calculation giving
the Schr\"odinger-picture density blocks and the support reduction producing
the zero summand of Wolf's Theorem~6.14~\cite{Wolf2012Quantum} are assembled
in the QICLean paper-gap note
\url{https://sirui-lu.com/QICLean/paper-gaps/wolf_theorem6_14_fixed_point_projection_gap.pdf}.

For the transported sector maps, the channel hypotheses must hold on the full
direct sums, not only on the contraction families. In the retained-coordinate
construction, compression by the target coisometry preserves trace exactly
when the precompression output belongs to its active image. The established
image identity is restricted to the source-generated contractions.

The direct-sum version of the dimension estimate is also complete. The proof
compresses to diagonal blocks, applies the direct-sum map, and embeds the
result block diagonally. Positivity, total-trace preservation, and the Schwarz
inequality pass to this extension. Its fixed points compress injectively into
the fixed points of the direct-sum map, while products of direct-sum fixed
points embed injectively among products of fixed points of the extension. The
density-block calculation for the extension therefore yields
\begin{align}
    \dim C_L(\mathcal F)
        &\leq \dim\mathcal F
    \label{eq:cpsv16_vertical_sector_invertibility_direct_sum_dimension_bound}
\end{align}
for the original direct sum.

Positivity, trace preservation, and the Schwarz hypothesis for both square
composites are resolved below. The canonical full-matrix lifts of the
transported maps factor through normalized state preparation,
retained-coordinate reindexing, the physical completely positive maps,
single-operator conjugation, and controlled partial trace. The transported
maps and their square composites are therefore completely positive. Once the
square composites are known to preserve trace, their trace adjoints are
unital completely positive maps and satisfy the Schwarz inequality in every
simple summand.\footnote{Formal declarations:
\leanid{MPOTensor.transportedVerticalSectorT_isKrausDirectSumMap},
\leanid{MPOTensor.transportedVerticalSectorS_isKrausDirectSumMap}, and
\leanid{MPOTensor.%
transportedVerticalSector_composites_traceAdjointSchwarz}
in \doclink{TNLean/MPS/MPDO/VerticalSectorCompletePositivity}.}

An alternative route uses the maximal-support witness
$T_\infty(\mathbf 1)$ for an arbitrary positive trace-preserving map. The
supported endomorphism and its complementary zero summand are constructed in
the full-matrix theorem recorded in the same QICLean paper-gap note. This
construction removes the need to assume a positive-definite fixed point before
applying the full-matrix classification.

\section{Resolution of the composite trace hypothesis}

The range inclusions
\begin{align}
    T(R_1(\mathcal A_1))
        &\subseteq R_2(\mathcal A_2),
    \label{eq:cpsv16_vertical_sector_invertibility_T_range_inclusion}\\
    S(R_2(\mathcal A_2))
        &\subseteq R_1(\mathcal A_1)
    \label{eq:cpsv16_vertical_sector_invertibility_S_range_inclusion}
\end{align}
have not been extracted from lines~1974--1980 of
Ref.~\cite{Cirac2016MPDO_arXiv}. They are unnecessary for trace preservation.
The argument first proves that both square composites preserve trace on their
full algebras and then derives the same property for each transported map by a
trace sandwich.

Let $F$ denote either square composite. The map $F$ is positive and trace
nonincreasing. Every forward orbit is therefore bounded: a positive orbit
remains positive with decreasing trace, and Hermitian decomposition gives the
general case. Let $P$ be the mean-ergodic projection and define
\begin{align}
    \rho
        &= P(\mathbf 1).
    \label{eq:cpsv16_vertical_sector_invertibility_ergodic_fixed_point}
\end{align}
The operator $\rho$ is positive and fixed, and its support contains every
fixed point of $F$. In particular, it contains every distinguished vertical
bond contraction. Since positive-length products of these contractions span
the whole sector algebra, their common support contains the identity.
Therefore $\rho$ is positive definite.

The corresponding full-matrix statement is formalized: a positive
trace-nonincreasing endomorphism whose fixed factors have products of one
positive length spanning the identity has a positive-definite fixed point. No
product is required to be fixed.\footnote{Formal declaration:
\leanid{IsPositiveMap.%
exists_posDef_fixedPoint_of_traceNonincreasing_of_fixed_product_span}
in \doclink{TNLean/Channel/FixedPoint/TraceNonincreasingProductSpan}.}
The finite-product version and its combination with the density-block
dimension bound are also formalized.\footnote{Formal declarations:
\leanid{Matrix.IsPositiveDirectSumMap.%
exists_posDef_fixedFamily_of_fixed_product_span}
in \doclink{TNLean/Channel/FixedPoint/TraceNonincreasingDirectSum}, and
\leanid{MPOTensor.%
eq_id_of_weightedVerticalBondContractions_fixed_of_traceAdjointSchwarz}
in \doclink{TNLean/MPS/MPDO/VerticalSectorInvertibility}.}

Thus a positive trace-preserving endomorphism of a finite product has a
positive-definite fixed family under the same product-span hypothesis. If its
trace adjoint satisfies the Schwarz inequality, the fixed contractions and
the density-block bound force the endomorphism to be the identity. For the
transported composites, complete positivity of the canonical lifts and trace
preservation supply the required trace-adjoint Schwarz inequalities from the
source tensor hypotheses. Applying the criterion to both fixed contraction
families gives
\begin{align}
    \widetilde S\widetilde T
        &= \operatorname{id}_{\mathcal A_1},
    \label{eq:cpsv16_vertical_sector_invertibility_first_composite_identity}\\
    \widetilde T\widetilde S
        &= \operatorname{id}_{\mathcal A_2}.
    \label{eq:cpsv16_vertical_sector_invertibility_second_composite_identity}
\end{align}
These are exactly the two composition identities in Appendix~C.4,
lines~1974--1995, of Ref.~\cite{Cirac2016MPDO_arXiv}.\footnote{Formal
declaration:
\leanid{MPOTensor.transportedVerticalSector_composites_eq_id}
in \doclink{TNLean/MPS/MPDO/VerticalSectorIdentity}.}

To prove trace preservation directly, let $F^*$ denote the trace adjoint and
define the trace-loss operator
\begin{align}
    \Delta
        &= \mathbf 1-F^*(\mathbf 1).
    \label{eq:cpsv16_vertical_sector_invertibility_trace_loss_operator}
\end{align}
For every positive $X$, trace nonincrease gives
\begin{align}
    \tr(\Delta X)
        &= \tr(X)-\tr(F(X)),
    \label{eq:cpsv16_vertical_sector_invertibility_trace_loss_identity}\\
    \tr(\Delta X)
        &\geq 0.
    \label{eq:cpsv16_vertical_sector_invertibility_trace_loss_nonnegative}
\end{align}
Positivity of $F$ makes $\Delta$ Hermitian, so
Eq.~\ref{eq:cpsv16_vertical_sector_invertibility_trace_loss_nonnegative}
implies $\Delta\geq0$. Since $\rho$ is fixed,
\begin{align}
    \tr(\rho\Delta)
        &= 0.
    \label{eq:cpsv16_vertical_sector_invertibility_fixed_trace_loss_zero}
\end{align}
Positive definiteness of $\rho$ then forces $\Delta=0$, which is equivalent
to trace preservation on the full algebra.

Now let $X\geq0$ belong to the one-site sector algebra. Positivity and trace
nonincrease give
\begin{align}
    \tr((\widetilde S\widetilde T)(X))
        &\leq \tr(\widetilde T(X)),
    \label{eq:cpsv16_vertical_sector_invertibility_trace_sandwich_first}\\
    \tr(\widetilde T(X))
        &\leq \tr(X).
    \label{eq:cpsv16_vertical_sector_invertibility_trace_sandwich_second}
\end{align}
The first and last quantities are equal by
Eq.~\ref{eq:cpsv16_vertical_sector_invertibility_first_composite_identity};
hence both inequalities are equalities. Thus $\widetilde T$ preserves trace
on the positive cone and, by linearity, on the whole algebra. The second
composite identity gives the same conclusion for $\widetilde S$.

This argument is formalized first for full matrix algebras, then for finite
direct sums through block-diagonal extension, and finally for the two
transported vertical-sector composites. It assumes neither multiplicativity
nor fixedness of products. Its only product hypothesis is the
basis-of-normal-tensors conclusion that products of the fixed contraction
factors span the ambient sector algebra. Trace preservation is therefore
closed for the individual maps and their square composites, although the
active-image inclusions in
Eqs.~\ref{eq:cpsv16_vertical_sector_invertibility_T_range_inclusion}
and~\ref{eq:cpsv16_vertical_sector_invertibility_S_range_inclusion} remain
unproved.

\section{Sector relabelling from the star-algebra equivalence}

After proving mutual invertibility, CPSV16 invokes the argument of
Wolf--P\'erez-Garc\'ia, Theorem~8
\cite[Appendix~C.4, lines~1995--2003]{Cirac2016MPDO_arXiv}. The cited theorem
concerns a trace-preserving Schwarz map. Its proof constructs a positive,
trace-preserving inverse on the peripheral space and uses the extreme points
of the density-operator sections
\begin{align}
    \bigoplus_k M_{d_k}(\mathbb C)\otimes\rho_k.
    \label{eq:cpsv16_vertical_sector_invertibility_density_operator_sections}
\end{align}
The map and its positive inverse send extreme density operators to extreme
density operators. Continuity and bijectivity then produce a permutation of
the simple summands and equality of the corresponding matrix dimensions. On
each matched summand, positivity in both directions permits either unitary
conjugation or transposed unitary conjugation; the Schwarz condition excludes
the transpose alternative
\cite[Theorem~8 and its proof]{WolfPerezGarcia2010Inverse}.

A linear equivalence between finite products of matrix spaces does not by
itself determine the simple summands. The trace adjoints of a mutually inverse
CPTP pair have been proved, by the Schwarz-defect argument, to be mutually
inverse $*$-algebra homomorphisms. Sector relabelling therefore reduces to two
algebraic steps:
\begin{enumerate}
    \item apply uniqueness of simple summands to the $*$-algebra equivalence,
          obtaining a sector permutation and equality of the paired matrix
          dimensions;
    \item identify each matched simple-summand isomorphism with unitary
          conjugation.
\end{enumerate}

The first step is proved for any $*$-algebra isomorphism between finite
products of nonzero full matrix algebras. The block ideals are exactly the
atoms of the two-sided ideal lattices, so the induced order isomorphism gives
an equivalence $\sigma$ between the block index sets. If $E$ denotes the
$*$-algebra isomorphism and $I_i,J_j$ are the source and target block ideals,
the defining matching relation is
\begin{align}
    E(I_i)
        &= J_{\sigma(i)}.
    \label{eq:cpsv16_vertical_sector_invertibility_block_ideal_matching}
\end{align}
The restriction to a paired block is a bijective complex-linear map
\begin{align}
    M_{d_i}(\mathbb C)
        &\longrightarrow M_{e_{\sigma(i)}}(\mathbb C).
    \label{eq:cpsv16_vertical_sector_invertibility_paired_block_map}
\end{align}
Comparison of dimensions gives
\begin{align}
    d_i
        &= e_{\sigma(i)}.
    \label{eq:cpsv16_vertical_sector_invertibility_paired_block_dimension}
\end{align}

Applying this result to the trace-adjoint $*$-algebra isomorphism proves
Eqs.~\ref{eq:cpsv16_vertical_sector_invertibility_block_ideal_matching}
and~\ref{eq:cpsv16_vertical_sector_invertibility_paired_block_dimension} for
every mutually inverse CPTP pair. The second step retains these witnesses. The
restriction to a paired block is a $*$-algebra isomorphism; after reindexing
with the retained dimension equality, Skolem--Noether gives an invertible
implementer $P$. Preservation of the adjoint forces $P^*P$ to be a positive
scalar matrix, and normalization gives a unitary implementer on the actual
target block.\footnote{Formal declarations:
\leanid{TwoSidedIdeal.exists_blockEquiv_of_ringEquiv_pi_simple} in
\doclink{TNLean/Algebra/BlockPermutation};
\leanid{Matrix.exists_unitary_conj_of_starAlgEquiv_matrix} and
\leanid{Matrix.%
exists_unitary_block_implementers_of_starAlgEquiv_pi_matrix} in
\doclink{TNLean/Algebra/SkolemNoetherUnitary};
\leanid{Matrix.exists_blockEquiv_dim_eq_of_starAlgEquiv_pi_matrix} and
\leanid{Matrix.%
exists_blockEquiv_dim_eq_unitary_of_mutual_inverse_kraus_direct_sum_maps} in
\doclink{TNLean/Channel/FixedPoint/DirectSumBlockPermutation}.}

The trace pairing for this unitary $*$-algebra isomorphism identifies it with
the original forward map, not only with the adjoint of the inverse. Its inverse
uses the same relabelling and the adjoints of the same unitaries. Applied to
the transported vertical-sector maps, this gives the two complete
single-summand formulas for $\widetilde T$ and $\widetilde S$ stated at
line~1997 of Ref.~\cite{Cirac2016MPDO_arXiv}.\footnote{Formal declarations:
\leanid{Matrix.%
exists_blockEquiv_dim_eq_unitary_forward_of_mutual_inverse_kraus_direct_sum_maps}
in \doclink{TNLean/Channel/FixedPoint/DirectSumBlockPermutation}, and
\leanid{MPOTensor.transportedVerticalSector_exists_unitaryBlockEquiv} in
\doclink{TNLean/MPS/MPDO/VerticalSectorRelabeling}, and
\leanid{MPOTensor.%
transportedVerticalSector_exists_unitaryBlockEquiv_coefficient_eq} in
\doclink{TNLean/MPS/MPDO/VerticalSectorCoefficientComparison}.}

\section{Elimination plan}

The first two steps of the Appendix~C.4 argument in
Ref.~\cite{Cirac2016MPDO_arXiv} are complete:
\begin{enumerate}
    \item complete positivity of the canonical lifts of the transported maps,
          and hence the adjoint Schwarz inequality for the two
          trace-preserving square composites, is established without an
          all-input active-image statement;
    \item the finite-product identity criterion is applied to each square
          composite using the fixed contraction families and the proved
          product-generation theorem, giving both identity compositions.
\end{enumerate}
The inverse-UCP multiplicative-domain step is also complete: the rectangular
trace adjoints form mutually inverse $*$-algebra equivalences. The generic
simple-summand correspondence and its internal unitary action give an
equivalence of the sector index sets, identify the block ideals carried into
one another, prove equality of paired matrix dimensions, and implement each
paired component by unitary conjugation. The specialization to both
transported maps retains the same relabelling and the same unitaries.

Finally, applying the transported refinement map to the weighted contraction
family gives
\begin{align}
    m_\alpha V_\alpha\iota_\alpha(M_\alpha(X))V_\alpha^*
        &= n_{\sigma(\alpha)}A_{\sigma(\alpha)}(X).
    \label{eq:cpsv16_vertical_sector_invertibility_weighted_contraction_comparison}
\end{align}
The positive multiplicity matrix $\nu_{\sigma(\alpha)}$ has nonzero trace, so
cancellation in
Eq.~\ref{eq:cpsv16_vertical_sector_invertibility_weighted_contraction_comparison}
gives
\begin{align}
    A_{\sigma(\alpha)}(X)
        &= \frac{m_\alpha}{n_{\sigma(\alpha)}}
           V_\alpha\iota_\alpha(M_\alpha(X))V_\alpha^*.
    \label{eq:cpsv16_vertical_sector_invertibility_trace_ratio_formula}
\end{align}
Evaluation on matrix units yields the tensor-letter identity stated at
line~2008 of Ref.~\cite{Cirac2016MPDO_arXiv}. This comparison concerns the
traces of the multiplicity matrices. It does not identify their dimensions,
the matrices themselves, or their individual diagonal entries.

\section{Classification}

\textbf{Verdict: the Appendix~C.4 argument through the coefficient identity at
line~2008 of Ref.~\cite{Cirac2016MPDO_arXiv} is closed.}
The transported maps and their square composites preserve trace under the
source tensor hypotheses. Both compositions are identities, and the
rectangular trace adjoints form mutually inverse $*$-algebra equivalences.
For every such pair, an equivalence of the index sets matches the simple
summands, the $*$-algebra isomorphism carries the corresponding block ideals
into one another, the paired matrix dimensions agree, and the induced map on
each matched summand is unitary conjugation. These formulas hold for
$\widetilde T$ and $\widetilde S$ themselves. Their action on the weighted
contraction family gives the trace-ratio coefficient formula in
Eq.~\ref{eq:cpsv16_vertical_sector_invertibility_trace_ratio_formula} and its
tensor-letter specialization.

The stronger active-image inclusions remain unproved and are unnecessary for
this route. No equality of multiplicities, multiplicity matrices, or
individual weights is asserted. No conclusion after Appendix~C.4, line~2008,
of Ref.~\cite{Cirac2016MPDO_arXiv} is included.

\bibliographystyle{alpha}
\bibliography{references}

\end{document}
